% ============================================================================
%  Subdocumento 1. Presentación de la empresa y organización del oferente
%  audIT Soluciones Tecnológicas SpA - Licitación TFEP-01/2026
% ============================================================================

\begin{resumenApertura}
audIT presenta sus capacidades institucionales y antecedentes relevantes para la Licitación Pública N.° TFEP-01/2026 de Transportes Curimón S.A. El documento resume la estructura y dotación de la empresa, su sistema de gestión de calidad y seguridad, la experiencia descrita en el Formulario T-6 y la capacidad de soporte declarada para la operación.
\begin{recibe}
  \item Estructura organizacional y asignación nominada de los seis roles habilitantes obligatorios del \art{34.1}{24}.
  \item Plataforma de borde embarcado audIT EdgeHub v2.4 con almacenamiento no volátil propuesto de 8~GB y objetivo de retención de al menos 72 horas, conforme al requisito \caso{RT-03.10}{31}; cualquier capacidad superior requiere validación y coordinación con la arquitectura del Subdocumento 4.
  \item Red regional de asistencia técnica y reemplazo de hardware en Ruta 5 con SLA contractual menor a cuatro horas.
  \item Certificación ISO 9001:2015 emitida y vigente; proceso ISO/IEC 27001:2022 en auditoría Fase 2 con Bureau Veritas y certificado comprometido para el Mes 1; y tres proyectos respaldados en T-6 con disponibilidades de 99,5\%, 99,2\% y 99,6\%.
\end{recibe}
\end{resumenApertura}

Este subdocumento presenta las capacidades técnicas, organizacionales y financieras declaradas por audIT para la Licitación Pública N.° TFEP-01/2026 de Transportes Curimón S.A. Resume la dotación de la empresa, su gobierno interno, la experiencia informada en el Formulario T-6 y las alianzas y capacidades de soporte pertinentes para el contrato.

El Subdocumento 1 complementa la comprensión del problema presentada en el Subdocumento 2 y resume las capacidades del oferente relevantes para el proyecto. El detalle de experiencia se presenta en el Formulario T-6 y los hitos de certificación, en el Anexo 1.A.

\section{Presentación de la empresa}\label{sec:1-presentacion}

audIT es una empresa chilena fundada en 2022 que presta servicios de ingeniería de software, arquitectura de sistemas distribuidos, telemetría y ciberseguridad operacional. Su actividad se orienta al diseño, desarrollo, despliegue y operación de plataformas digitales para organizaciones logísticas con infraestructura geográficamente dispersa y conectividad intermitente.

A continuación se detallan la trayectoria, principios de gestión, líneas de servicio formalizadas, catálogo de productos tecnológicos y la infraestructura logística y territorial de la compañía.

\subsection{Trayectoria institucional, misión y presencia geográfica}\label{sec:1-trayectoria}

Desde su constitución, audIT ha concentrado su propuesta de valor en resolver la brecha de continuidad operacional que experimentan los sistemas convencionales cuando deben sincronizar flujos de datos capturados en terreno con plataformas centrales en la nube. La empresa dispone de infraestructura propia de investigación, desarrollo y control de calidad distribuida en dos centros operacionales:

\begin{enumerate}
  \item \textbf{Casa Matriz y Centro de Ingeniería de Software:} Ubicado en la ciudad de Viña del Mar (Región de Valparaíso), donde residen las células de arquitectura cloud, desarrollo backend, ingeniería de datos y la oficina de gestión de proyectos (PMO).
  \item \textbf{Centro de Operaciones Técnicas, Laboratorio de Hardware y Banco de Pruebas Telemáticas:} Ubicado en la comuna de Santiago (Región Metropolitana), equipado con instrumental de laboratorio electrónico, bancos de simulación de bus CAN automotriz bajo estándar SAE J1939, cámaras térmicas de prueba ambiental y estaciones de precarga de firmware seguro. Esta ubicación intermedia asegura un despliegue ágil sobre el eje logístico central Valparaíso -- Santiago -- San Bernardo -- Los Andes.
\end{enumerate}

Los principios fundacionales que gobiernan el accionar corporativo de audIT se definen como sigue:

\begin{itemize}
  \item \textbf{Giro Corporativo Principal:} Servicios integrales de consultoría informática, desarrollo de software a medida, diseño e implantación de arquitecturas cloud híbridas, ingeniería de integración telemática e Internet de las Cosas (IoT) industrial, y auditoría de ciberseguridad y protección de datos bajo estándares internacionales.
  \item \textbf{Misión Corporativa:} Proveer soluciones de ingeniería de software y telemetría distribuida de alta disponibilidad que garanticen la continuidad operativa, la trazabilidad auditable y la seguridad de la información en entornos logísticos e industriales de misión crítica, sustentadas en evidencia técnica rigurosa y cumplimiento normativo estricto.
  \item \textbf{Visión Corporativa:} Desarrollar servicios de ingeniería para la transformación digital de cadenas de suministro y operaciones de transporte, incluidas plataformas de datos tolerantes a desconexión (\textit{offline-first}).
  \item \textbf{Valores de Ingeniería:} Rigor metodológico, seguridad por diseño (\textit{Zero Trust}), transparencia auditable y cumplimiento de los acuerdos de nivel de servicio (SLA) convenidos con los mandantes.
\end{itemize}

\subsection{Líneas de servicio de ingeniería especializada}\label{sec:1-lineas-servicio}

audIT estructura sus capacidades instaladas en tres líneas de servicio de ingeniería formal, cada una dotada de procedimientos estandarizados, metodologías certificadas, entregables medibles y marcos de responsabilidad técnica definidos:

\begin{enumerate}
  \item \textbf{Línea de Ingeniería de Telemetría de Borde y Sistemas Ciber-Físicos:}
  \begin{itemize}
    \item \textit{Alcance y Responsabilidad Técnica:} Diseño, programación, integración y mantenimiento de firmware industrial para dispositivos electrónicos embarcados en cabina y nodos de campo. Integración de telemetría vehicular multimarca mediante lectura pasiva no intrusiva de buses CAN automotrices bajo especificaciones SAE J1939, FMS (\textit{Fleet Management System}) y tacógrafos digitales, asegurando la no alteración del cableado original ni la pérdida de garantías de fabricante de los vehículos. Sensorización especializada para cadena de frío (monitoreo de temperatura multizona con sondas PT100) y control de accionamiento de tomas de fuerza y válvulas en tolvas y estanques.
    \item \textit{Entregables Formales:} Paquetes de firmware firmados criptográficamente, pasarelas telemáticas preconfiguradas, protocolos de prueba de compatibilidad electromagnética y eléctrica en cabina, y manuales de homologación de lectura de buses de datos.
  \end{itemize}

  \item \textbf{Línea de Plataformas Cloud Resilientes e Ingeniería de Datos:}
  \begin{itemize}
    \item \textit{Alcance y Responsabilidad Técnica:} Arquitectura, desarrollo, orquestación y administración continuada de plataformas analíticas y transaccionales en la nube bajo esquemas híbridos y multirregión. Implementación de capas de ingestión masiva de eventos en tiempo real con capacidad de absorción de ráfagas, bases de datos optimizadas para series temporales y modelos relacionales, motores algorítmicos de optimización de rutas y despacho, e interfaces de usuario web y móviles de alto rendimiento y diseño ergonómico. Implementación de Capas Anticorrupción (\textit{Domain-Driven Design}) para integrar e interoperar de manera transparente con sistemas heredados (\textit{legacy}) y plataformas telemáticas externas de terceros sin degradar la consistencia de los datos.
    \item \textit{Entregables Formales:} Repositorios de código fuente bajo control de versiones, catálogos de interfaces de programación de aplicaciones (APIs RESTful seguras bajo estándar OpenAPI 3.1), esquemas y modelos de datos normalizados, tuberías automatizadas de despliegue (CI/CD) y tableros analíticos e informes ejecutivos de rendimiento operacional.
  \end{itemize}

  \item \textbf{Línea de Ciberseguridad Operacional, Gobernanza y Continuidad de Negocio:}
  \begin{itemize}
    \item \textit{Alcance y Responsabilidad Técnica:} Diseño de controles de seguridad de la información tomando como referencias ISO/IEC 27001:2022 y NIST CSF 2.0. Se consideran Arquitectura de Confianza Cero (\textit{Zero Trust}, NIST SP 800-207), cifrado de datos en tránsito y en reposo, controles de acceso y protección de datos personales. La Ley N.° 21.719 fue promulgada, pero su entrada en vigor está fijada para el 1 de diciembre de 2026; para la fecha de esta oferta, el marco vigente debe revisarse junto con la Ley N.° 19.628. El diseño deberá actualizarse y validarse antes de la puesta en servicio. También se consideran prácticas de cadena de suministro de software, OWASP ASVS y continuidad operacional, sujetas a definición de alcance y evidencia (\caso{RT-11.28}{33}).
    \item \textit{Entregables Formales:} Informes de pruebas de penetración y análisis estático/dinámico de vulnerabilidades (SAST/DAST), matrices de riesgos de seguridad, registros contables criptográficos inmutables de auditoría (\textit{WORM audit logs}), políticas formales de privacidad y planes de continuidad de negocio auditados.
  \end{itemize}
\end{enumerate}

\subsection{Plataforma de borde embarcado: audIT EdgeHub v2.4 Enterprise}\label{sec:1-edgehub}

En el ámbito de productos propios, audIT ha desarrollado y estandarizado la solución \textbf{audIT EdgeHub v2.4 Enterprise Embedded}, un paquete integral de software de borde (\textit{edge computing}) concebido específicamente para operar en pasarelas telemáticas instaladas a bordo de tractocamiones y vehículos de carga pesada que transitan por corredores con cobertura de telecomunicaciones intermitente o nula.

Las especificaciones técnicas y operacionales del producto audIT EdgeHub v2.4 Enterprise se resumen a continuación:

\begin{itemize}
  \item \textbf{Entorno de Ejecución y Sistema Operativo:} Basado en distribución Linux industrial embebida (\textit{Debian Embedded / Yocto Project}), con núcleo (\textit{kernel}) endurecido conforme a guías CIS Benchmarks, arranque seguro (\textit{Secure Boot}) y particionamiento de almacenamiento redundante A/B para permitir actualizaciones remotas de firmware (\textit{FOTA -- Firmware Over-The-Air}) a prueba de fallos y cortes intempestivos de energía.
  \item \textbf{Lenguajes y módulos de control:} La propuesta considera módulos de captura cinemática y comunicación con el bus CAN en Rust y C/C++. El consumo de recursos, la seguridad de memoria y la recuperación ante interrupciones deberán comprobarse mediante pruebas sobre el hardware seleccionado.
  \item \textbf{Búfer local y resiliencia eléctrica:} La propuesta considera SQLite con registro \textit{Write-Ahead Logging} (WAL) en memoria eMMC industrial de al menos 8~GB (\caso{RT-08.11}{32}). WAL mantiene un registro de transacciones que puede facilitar la recuperación tras una interrupción, pero no sustituye la configuración de sincronización, el manejo de energía ni las pruebas de integridad. La recuperación y la retención de datos deben validarse con ensayos de corte eléctrico y reinicio en el dispositivo seleccionado; no se declara ausencia garantizada de pérdida o corrupción sin esos resultados.
  \item \textbf{Autonomía operacional y sincronización (RT-03.10):} El búfer local propuesto usa memoria eMMC industrial de 8~GB y debe retener al menos 72 horas de datos en zonas sin cobertura. La capacidad debe verificarse con el perfil efectivo de muestreo, la carga de eventos, el espacio disponible después del sistema operativo y las pruebas de desconexión. El Subdocumento 4 podrá proponer una capacidad superior, sujeta a cálculo y validación en banco.
  \item \textbf{Seguridad en el borde y en cabina:} Se propone evaluar un elemento seguro para la custodia de claves y la firma de eventos, además de cifrado de datos sensibles. Los controles de privacidad deberán alinearse con la normativa vigente y prepararse para la entrada en vigor de la Ley N.° 21.719 el 1 de diciembre de 2026. Para la interacción en cabina, el diseño considera reducir distracciones y restringir operaciones que puedan apartar la atención del conductor; el alcance debe cotejarse con la Ley N.° 21.377 y su reglamentación antes de definir bloqueos o alertas de voz como requisito.
\end{itemize}

\subsection{Red regional de asistencia técnica y reemplazo de hardware en Ruta 5}\label{sec:1-red-regional}

Para asegurar una respuesta operativa expedita frente a incidencias físicas de hardware en terreno y garantizar el cumplimiento irrestricto de las exigencias establecidas en el Capítulo 8.4 y en el requerimiento \transv{RT-21.16}{36} de las Bases Técnicas Transversales, audIT ha articulado y formalizado una \textbf{Red de Soporte Regional en Terreno} mediante \textbf{Convenios Marco de Prestación de Servicios de Soporte y Acuerdos de Nivel Operacional (OLA)} legalmente vinculantes con cuatro centros técnicos y maestranzas automotrices especializadas, estratégicamente distribuidas a lo largo de los 3.000 kilómetros del corredor de la Ruta 5:

\begin{enumerate}
  \item \textbf{Nodo Regional Antofagasta (Macrozona Norte):}
  \begin{itemize}
    \item \textit{Centro Técnico Autorizado:} Maestranza y Laboratorio Electrónico Diesel del Norte Ltda. · RUT: 76.512.890-4.
    \item \textit{Instrumento Jurídico:} Convenio Marco de Soporte y OLA N.° CM-2025-01-ANF.
    \item \textit{Emplazamiento Físico:} Avenida Pedro Aguirre Cerda N.° 8450, Barrio Industrial La Chimba, Antofagasta.
    \item \textit{Radio de Cobertura y Función Operativa:} Cobertura prioritaria para las operaciones mineras, hubs de sustancias peligrosas (D.S. 298) y faenas portuarias del Norte Grande, proveyendo asistencia técnica directa en el Terminal Antofagasta de Transportes Curimón S.A.
  \end{itemize}

  \item \textbf{Nodo Regional Talca (Macrozona Centro-Sur):}
  \begin{itemize}
    \item \textit{Centro Autorizado:} Electromecánica y Telecomunicaciones del Maule SpA (RUT: 77.104.530-K).
    \item \textit{Instrumento Jurídico:} Convenio Marco de Soporte y OLA N.° CM-2025-02-TAL.
    \item \textit{Emplazamiento Físico:} Longitudinal Sur Km 252, Cruce Varoli, Talca.
    \item \textit{Radio de Cobertura y Función Operativa:} Cobertura estratégica del corredor agroindustrial y de transferencia de carga entre Rancagua y Chillán, asegurando atención inmediata sobre la Ruta 5 Sur.
  \end{itemize}

  \item \textbf{Nodo Regional Los Ángeles (Macrozona Forestal / Biobío):}
  \begin{itemize}
    \item \textit{Centro Técnico Autorizado:} Centro Integral de Servicios Telemáticos Biobío S.A. · RUT: 76.890.120-3.
    \item \textit{Instrumento Jurídico:} Convenio Marco de Soporte y OLA N.° CM-2025-03-LAN.
    \item \textit{Emplazamiento Físico:} Avenida Las Industrias N.° 5200, Ruta 5 Sur Km 512, Los Ángeles.
    \item \textit{Radio de Cobertura y Función Operativa:} Soporte dedicado para operaciones forestales, industriales y rutas transversales hacia Concepción, Talcahuano y la Cordillera de Nahuelbuta.
  \end{itemize}

  \item \textbf{Nodo Regional Puerto Montt (Macrozona Sur / Austral):}
  \begin{itemize}
    \item \textit{Centro Técnico Autorizado:} Laboratorio y Soporte Tecnológico Austral Ltda. · RUT: 76.621.904-7.
    \item \textit{Instrumento Jurídico:} Convenio Marco de Soporte y OLA N.° CM-2025-04-PMC.
    \item \textit{Emplazamiento Físico:} Ruta 5 Sur Km 1025, Sector Alto Cardonal, Puerto Montt.
    \item \textit{Radio de Cobertura y Función Operativa:} Cobertura de la industria acuícola, centros de cultivo de salmón y enlace con la zona austral y el Terminal Puerto Montt de Transportes Curimón S.A.
  \end{itemize}
\end{enumerate}

La exigibilidad de estos acuerdos garantiza:
\begin{itemize}
  \item \textbf{Compromiso de Reemplazo (SLA $< 4\text{ h}$):} Los convenios suscritos obligan contractualmente a los aliados a mantener cuadrillas técnicas de turno rotativo las 24 horas del día, los 365 días del año, garantizando un tiempo máximo de restitución o reemplazo físico de pasarelas telemáticas, sensores térmicos o cableado dañado inferior a cuatro (4) horas corridas contadas desde el despacho de la orden de servicio.
  \item \textbf{Stock Crítico de Repuestos en Custodia Legal:} En virtud de los Acuerdos OLA, cada centro autorizado mantiene en custodia un inventario permanente de seguridad de propiedad de audIT, auditado mensualmente y trazable por número de serie: 10 pasarelas vehiculares \textit{audIT EdgeHub v2.4} completas y preconfiguradas; 15 kits de acopladores inductivos \textit{CANclick}; 15 sondas de temperatura PT100 con cable siliconado industrial; 10 antenas combinadas GNSS/4G-LTE IP67; y cableado automotriz retardante a la llama con fusibles de protección vehicular.
\end{itemize}

\subsection{Cumplimiento copulativo de la normativa de sustancias peligrosas}\label{sec:1-suspel}

Dentro de la flota operada por Transportes Curimón S.A. se contabilizan \textbf{18 unidades de transporte especializadas en sustancias peligrosas (SUSPEL)}. audIT comprende a cabalidad que el marco regulatorio chileno impone exigencias diferenciadas y copulativas para este segmento de carga crítica, distinguiendo formalmente dos ámbitos normativos que deben cumplirse de manera conjunta:

\begin{itemize}
  \item \textbf{Decreto Supremo N.° 298/1994 del Ministerio de Transportes y Telecomunicaciones:} Rige estrictamente el \textbf{transporte de cargas peligrosas por calles, caminos y carreteras públicas}. audIT asegura la conformidad operacional de las 18 unidades en ruta mediante la monitorización satelital continua de velocidad, trazabilidad estricta de rutas autorizadas, control de paradas en sitios no habilitados, geocercas dinámicas de advertencia ante proximidad a centros poblados o fuentes de agua, y verificación de la operatividad ininterrumpida de los dispositivos de registro telemático a bordo.
  \item \textbf{Decreto Supremo N.° 43/2015 del Ministerio de Salud:} Aprueba el reglamento de \textbf{almacenamiento de sustancias peligrosas}, aplicable con carácter copulativo en las instalaciones fijas, terminales de transferencia, recintos industriales y patios químicos donde las 18 unidades de carga peligrosa inician viaje, efectúan faenas de trasvasije o carguío, pernoctan en detención programada o concluyen su ciclo logístico. La plataforma de audIT integra el registro y control de permanencia temporal en patios de estanqueidad, verificación del cumplimiento de capacidades de contención secundaria, control de incompatibilidades químicas de adyacencia de carga en patios de maniobra y generación automática de bitácoras de fiscalización exigidas por la autoridad sanitaria.
\end{itemize}

\section{Estructura organizacional}\label{sec:1-organizacion}

El modelo de gestión institucional de audIT se estructura bajo un enfoque matricial de ingeniería de software y sistemas ciber-físicos, diseñado para otorgar agilidad en el desarrollo de software y rigor estricto en la homologación de hardware y soporte en terreno.

A continuación se presenta el organigrama corporativo, seguido de un análisis pormenorizado de las unidades funcionales, la dotación de 22 profesionales de planta y el esquema de asignación a los roles habilitantes exigidos por las Bases Administrativas.

\subsection{Organigrama institucional y dotación de planta}\label{sec:1-organigrama}

La estructura funcional de audIT se encuentra dimensionada para asegurar gobernanza centralizada, separación de responsabilidades y ejecución técnica pareada. En la \figref{1-organigrama} se ilustra la articulación orgánica de la compañía y sus líneas de mando.

\figura{figuras/01-empresa/Estructura_y_Organizacion.png}{Organigrama institucional y estructura operativa de audIT SpA}{1-organigrama}

Como se aprecia en la \figref{1-organigrama}, la dotación permanente de la compañía está constituida por exactamente \textbf{22 profesionales de planta} contratados bajo régimen laboral indefinido, cuya distribución analítica por unidad organizativa corresponde a la siguiente:

\begin{enumerate}
  \item \textbf{Gerencia General / Dirección de Proyectos (1 profesional):} Ingeniero Civil Informático con postgrado en gestión de proyectos y certificación PMP®. Concentra la representación legal de la sociedad, la dirección estratégica, el control presupuestario y la interlocución de máximo nivel con el mandante.
  \item \textbf{Oficina de Gestión de Proyectos (PMO) y Soporte Técnico (3 profesionales):} Tres analistas de ingeniería industrial e informática con certificaciones ITIL 4 Foundation. Responsables de la administración de la planificación contractual, control de la EDT, gestión de adquisiciones de hardware, seguimiento de hitos de la Carta Gantt y supervisión de la mesa de ayuda de soporte operacional de segundo nivel.
  \item \textbf{Área de Aseguramiento de Calidad, Ciberseguridad y Auditoría ISO (4 profesionales):} Cuatro especialistas de dedicación transversal e independientes de las células de desarrollo:
  \begin{itemize}
    \item \textit{1 Líder de Calidad de Software:} Ingeniero Civil Informático, con certificación ISTQB (\textit{Full} \textit{Advanced} \textit{Level}), responsable del plan de pruebas, calidad del código y Quality Gates.
    \item \textit{1 Encargado de Seguridad de la Información (CISO):} Especialista en ciberseguridad industrial, certificado CISSP y CISM, responsable del SGSI, criptografía y la planificación de controles de privacidad conforme a la normativa vigente y a la futura entrada en vigor de la Ley N.° 21.719.
    \item \textit{1 Auditor Interno de Procesos y Gobierno TI:} Ingeniero de Sistemas, certificado COBIT y Auditor Líder ISO 9001 / ISO 27001, a cargo de los ciclos de auditoría interna y trazabilidad normativa.
    \item \textit{1 Ingeniero de Pruebas y Automatización QA:} Ingeniero de Software dedicado al desarrollo de scripts de pruebas automatizadas, pruebas de carga y estrés en clúster cloud.
  \end{itemize}
  \item \textbf{Dirección de Arquitectura y Datos (1 profesional):} Ingeniero Civil Informático, Máster en Ciencias de la Computación, con certificaciones \textit{Azure Solutions Architect Expert} y \textit{CDMP}. Conduce la gobernanza tecnológica, los patrones de integración de microservicios, el diseño de la capa anticorrupción y el modelamiento de datos.
  \item \textbf{Células de Software, Backend y Plataforma Cloud (7 ingenieros):} Equipo multidisciplinario conformado por 3 ingenieros de software senior (microservicios, .NET Core, Go y Rust), 2 ingenieros de frontend y aplicaciones móviles (React, TypeScript y arquitecturas PWA), y 2 ingenieros de datos y bases de datos cloud (PostgreSQL, TimescaleDB, Azure Event Hubs y Kafka).
  \item \textbf{Jefatura de IoT y Sistemas de Terreno (1 profesional):} Ingeniero en Telecomunicaciones y Electrónica con más de una década de experiencia práctica en diseño de circuitos automotrices, integración de buses vehiculares SAE J1939, normativas FMS y protocolos de radiocomunicación celular.
  \item \textbf{Células de Firmware, Telemetría y Laboratorio de Hardware (5 ingenieros):} Equipo compuesto por 3 ingenieros de firmware embebido (C/C++ y Rust sobre Linux industrial embebido), 1 ingeniero de diseño electrónico y compatibilidad electromagnética (bancos de prueba CAN y sensorización de cabina), y 1 ingeniero de telecomunicaciones e integración de campo, responsable técnico de la coordinación con la red de centros autorizados regionales en la Ruta 5.
\end{enumerate}

\subsection{Roles habilitantes y asignación de responsabilidades}\label{sec:1-roles-habilitantes}

En conformidad con lo prescrito en el \art{34.1}{24} de las Bases Administrativas, audIT asigna e individualiza de manera unívoca a sus profesionales de planta más calificados para cubrir la totalidad de los seis (6) Roles Habilitantes obligatorios del contrato:

\begin{enumerate}
  \item \textbf{Jefe de Proyecto:} \textbf{Alejandro Hermosilla}. Ingeniero Civil Informático, certificado PMP®, con 14 años de experiencia en proyectos de modernización de plataformas digitales para logística y transporte.
  \item \textbf{Arquitecto de Solución:} \textbf{Dr. Esteban Valenzuela}. Ingeniero Civil Informático, con estudios de posgrado en Ciencias de la Computación y certificaciones \textit{Microsoft Certified: Azure Solutions Architect Expert} y \textit{TOGAF 9.2 Certified}; cuenta con 12 años de experiencia en arquitecturas cloud híbridas y microservicios.
  \item \textbf{Oficial de Seguridad de la Información (CISO):} \textbf{Mauricio Arancibia}. Especialista en ciberseguridad industrial, con certificaciones CISSP, CISM y Auditor Líder ISO/IEC 27001, y 11 años de experiencia en criptografía y gobernanza de protección de datos.
  \item \textbf{Líder de Datos:} \textbf{Rodrigo Sanhueza Parra}. Ingeniero Civil en Informática, certificado CDMP y \textit{Azure Data Engineer Associate}, con 9 años de experiencia en arquitectura de datos distribuidos, pipelines de eventos en tiempo real (Kafka / Azure Event Hubs) y bases de datos relacionales y de series temporales (TimescaleDB / PostgreSQL).
  \item \textbf{Líder de Aseguramiento de Calidad (QA):} \textbf{Claudia Navarrete Rivas}. Ingeniera Civil Informática, certificada ISTQB (\textit{Full Advanced Level} -- Test Manager, Technical Test Analyst), con 10 años de experiencia dirigiendo compuertas de calidad (\textit{Quality Gates}), pruebas automatizadas de regresión, carga y estrés, y aseguramiento normativo bajo ISO/IEC/IEEE 12207.
  \item \textbf{Líder de Operaciones / Infraestructura (SRE / DevOps):} \textbf{Felipe Morales Cárdenas}. Ingeniero en Telecomunicaciones, Conectividad y Redes, certificado CKA (Cloud Native Computing Foundation) y \textit{Red Hat Certified Engineer}, con 10 años de experiencia en orquestación de clústeres Kubernetes (AKS), telemetría industrial IoT, redes vehiculares SAE J1939 y soporte operacional 24/7/365 en terreno.
\end{enumerate}

La tabla anterior identifica los seis roles habilitantes del \art{34.1}{24}. Los antecedentes curriculares, certificados y cartas se encuentran respaldados en el Formulario T-8; su versión entregable debe cotejarse junto con el Formulario T-22 antes del cierre.

\subsection{Modelo operativo en células y mitigación de dependencia}\label{sec:1-celulas}

audIT opera mediante una metodología de ingeniería colaborativa en parejas (\textit{pair\allowbreak-engineering}) y células multifuncionales. Esta disciplina de diseño previene la generación de silos de conocimiento técnico y elimina el riesgo de dependencia de personas únicas (\textit{single} \textit{points} \textit{of} \textit{knowledge} \textit{failure}). Todo desarrollo de firmware, código backend o configuración de infraestructura en la nube es sometido a revisión cruzada obligatoria (\textit{peer review}) mediante solicitudes de integración (\textit{pull requests}) antes de incorporarse a la rama principal de compilación, resguardando la transferibilidad inmediata de funciones ante eventuales contingencias o reemplazos de personal.

\section{Gobierno interno Calidad, Seguridad y Conocimiento}\label{sec:1-gobierno}

El gobierno institucional de audIT articula de manera sinérgica la calidad del ciclo de vida del software, la seguridad de la información y la gestión del conocimiento técnico. Este marco garantiza que cada producto de trabajo cumpla con estándares auditables antes de ser transferido al entorno operacional del mandante.

A continuación se describen las instancias colegiadas de gobierno, los ciclos anuales de auditoría bajo normas internacionales, la matriz de escalamiento técnico y la estrategia corporativa de gestión del conocimiento.

\subsection{Comités institucionales de gobierno y cadencias de decisión}\label{sec:1-comites}

audIT formaliza su gobernanza mediante tres comités colegiados permanentes con calendarios y atribuciones estrictas:

\begin{enumerate}
  \item \textbf{Comité de Calidad y Procesos de Ingeniería (CCPI):}
  \begin{itemize}
    \item \textit{Periodicidad:} Sesiones ordinarias mensuales; sesiones extraordinarias ante desviaciones de calidad críticas.
    \item \textit{Composición:} Líder de Calidad (Presidente), Director de Proyectos, Director de Arquitectura y Datos, y Líder de Operación.
    \item \textit{Atribuciones y Responsabilidades:} Evaluación mensual de métricas de proceso (cobertura de pruebas, tasa de reapertura de defectos, cumplimiento de estándares de código limpio), aprobación formal de las compuertas de calidad (\textit{Quality Gates}) para transiciones de fase contractual, y autorización de pases a producción.
  \end{itemize}

  \item \textbf{Comité de Seguridad de la Información y Ciberseguridad (CSIC):}
  \begin{itemize}
    \item \textit{Periodicidad:} Sesiones ordinarias mensuales.
    \item \textit{Composición:} Encargado de Seguridad de la Información (Presidente), Gerente General, Director de Arquitectura y Auditor Interno de Procesos.
    \item \textit{Atribuciones y Responsabilidades:} Monitoreo y actualización de la matriz de riesgos bajo estándar NIST CSF 2.0 e ISO 27001:2022; revisión mensual de los registros inmutables de acceso a datos sensibles de conductores; auditoría de claves criptográficas y módulos HSM; y análisis de alertas de ciberseguridad emitidas por el CSIRT Nacional conforme a la Ley N.° 21.663.
  \end{itemize}

  \item \textbf{Comité Ejecutivo de Gobernanza y Continuidad de Negocio (CEGCN):}
  \begin{itemize}
    \item \textit{Periodicidad:} Sesiones trimestrales.
    \item \textit{Composición:} Gerente General, Director de Proyectos, Jefaturas de Área y Representante Legal.
    \item \textit{Atribuciones y Responsabilidades:} Revisión estratégica del desempeño de los contratos de largo plazo, evaluación de indicadores de solvencia y liquidez, análisis de resultados de auditorías externas ISO y validación de planes de contingencia operacional ante desastres (ISO 22301).
  \end{itemize}
\end{enumerate}

\subsection{Ciclo anual de auditorías y gestión de mejora continua}\label{sec:1-auditorias}

La compañía ejecuta un riguroso ciclo anual de aseguramiento de calidad y seguridad de la información, estructurado en cuatro hitos programados:

\begin{itemize}
  \item \textbf{Auditorías Internas Semestrales de Calidad (ISO 9001:2015):} En los meses de mayo y noviembre de cada año, el equipo de auditoría interna evalúa el 100\% de los procesos de desarrollo de software, pruebas de banco de hardware y gestión de mesa de ayuda, verificando la trazabilidad de requisitos y la aplicación de los planes de pruebas.
  \item \textbf{Auditorías Internas Semestrales de Seguridad (ISO/IEC 27001:2022):} En los meses de marzo y septiembre, se audita el cumplimiento de los 93 controles del Anexo A de la norma, con énfasis en el control de acceso, gestión de incidentes, cifrado en reposo y tránsito, y segregación lógica de entornos.
  \item \textbf{Auditoría Externa Anual de Mantenimiento / Recertificación:} Ejecutada por un organismo certificador internacional independiente debidamente acreditado (como Bureau Veritas o SGS), validando formalmente la eficacia del sistema integrado de gestión.
  \item \textbf{Procedimiento de Acciones Correctivas y Preventivas (CAPA):} Toda no conformidad detectada en auditorías internas o externas se incorpora de inmediato a un flujo de trabajo formal en el sistema de gestión corporativo, exigiendo análisis de causa raíz mediante metodología de los 5 Porqués o diagrama de Ishikawa, con un plazo perentorio de cierre y verificación de eficacia no superior a treinta (30) días calendario.
\end{itemize}

\subsection{Matriz institucional de escalamiento técnico y operacional}\label{sec:1-escalamiento}

Para garantizar tiempos de respuesta deterministas ante incidentes técnicos en las plataformas y hardware en ruta, audIT aplica una matriz de escalamiento estructurada en cuatro niveles jerárquicos:

\begin{itemize}
  \item \textbf{Nivel 1 (Atención Inicial y Diagnóstico de Mesa de Ayuda):} Analistas de Soporte Técnico (PMO y Soporte). Recepción de alertas automatizadas de telemetría y tickets de usuarios; diagnóstico inicial mediante consolas de monitoreo; aplicación de procedimientos estándar de operación (SOP). Tiempo máximo de escalamiento: 15 minutos ante incidentes clasificados como Críticos (Severidad 1).
  \item \textbf{Nivel 2 (Ingeniería de Especialidad y Soporte de Campo):} Ingenieros de Células de Software o Células de Firmware/Hardware. Depuración de código, análisis de registros de logs en Azure, emisión de parches de software urgentes o movilización del centro técnico regional en la Ruta 5 para reemplazo físico de hardware. Tiempo máximo de escalamiento: 45 minutos si la incidencia compromete la continuidad operacional del mandante.
  \item \textbf{Nivel 3 (Dirección Técnica de Especialidad):} Director de Arquitectura y Datos / Jefatura de IoT y Terreno / CISO. Decisión de activación de contingencia, conmutación por error (\textit{failover}) a región secundaria en nube, modificación de esquemas de red o instrucción de procedimientos de aislamiento de seguridad ante sospecha de intrusión. Tiempo máximo de escalamiento: 60 minutos.
  \item \textbf{Nivel 4 (Comité de Crisis y Dirección Ejecutiva):} Gerente General / Director de Proyectos y Contraparte Gerencial de Transportes Curimón S.A. Notificación formal a la alta dirección del mandante, activación del Plan de Continuidad de Negocio (BCP) y coordinación legal/comunicacional si correspondiera.
\end{itemize}

\subsection{Gestión del conocimiento y repositorios institucionales}\label{sec:1-conocimiento}

audIT gestiona su capital intelectual como un activo corporativo auditable a través de las siguientes herramientas y protocolos:

\begin{itemize}
  \item \textbf{Repositorio Central de Arquitecturas y Patrones:} Wiki técnica centralizada y versionada donde se documentan los diagramas de arquitectura de referencia, guías de estilo de código, protocolos de integración con buses CAN vehiculares y manuales de empaquetado de firmware.
  \item \textbf{Bitácora de Incidentes y Post-Mortems sin Culpa (\textit{Blameless Post-Mortems}):} Cada incidente de severidad alta genera obligatoriamente un informe técnico retrospectivo que desglosa la cronología del evento, el impacto en usuarios, la causa raíz tecnológica y las medidas preventivas implementadas en el código o en la infraestructura para evitar su recurrencia.
  \item \textbf{Programa Continuo de Entrenamiento y Certificación:} Plan anual corporativo de perfeccionamiento que financia la certificación técnica del personal de planta en tecnologías de nube Azure, ciberseguridad, gestión de datos y pruebas de software.
\end{itemize}

\section{Experiencia y certificaciones}\label{sec:1-experiencia}

La experiencia declarada por audIT se resume mediante proyectos del sector logístico y de transporte y certificaciones institucionales. Los antecedentes financieros corporativos, incluidos ratios y respaldos, se presentan en el Sobre N.º 1. El Formulario T-6 contiene el detalle de los proyectos y su documentación de respaldo.

La siguiente sección resume los proyectos similares y las certificaciones declaradas; no reproduce información financiera del Sobre N.º 1.

\subsection{Resumen analítico de proyectos homólogos}\label{sec:1-proyectos-homologos}

En conformidad con lo dispuesto en el \art{34}{24} de las Bases Administrativas y el Comunicado 10, en la \tabref{1-proyectos} se sintetizan los tres proyectos de misión crítica desarrollados y concluidos exitosamente por audIT dentro de los últimos cinco años. Los tres proyectos operan en la actualidad bajo esquemas de arquitectura híbrida, con disponibilidades mensuales de 99,5\%, 99,2\% y 99,6\%, y volúmenes operacionales comparables con la flota (374 camiones) y viajes anuales ($\approx 96.000$) de Transportes Curimón S.A.

Las disponibilidades que siguen son antecedentes históricos de los proyectos del T-6. Para el servicio crítico ofertado a Curimón, el nivel contractual aplicable es al menos 99,9\% de disponibilidad mensual (\art{20}{14}; \transv{RT-10.01}{22}).

\begin{tabla}[ancho=completo]{Resumen analítico de proyectos de experiencia previa acreditados}{1-proyectos}{L{34mm} X X X L{30mm}}
\textbf{Parámetro Clave} & \textbf{Proyecto 1: Transporte Interurbano} & \textbf{Proyecto 2: Logística en Frío y Faenas} & \textbf{Proyecto 3: Distribución Multimodal} & \textbf{Exigencia Art. 34° / Curimón} \\ \midrule
Cliente y Giro & Transportes del Sur Ltda. (Carga pesada interurbana) & AgroFrutícola Los Andes S.A. (Logística en frío remota) & Logística Multimodal Bicentenario S.A. (Carga portuaria) & Industria del caso y afines (Transporte y Logística) \\
Volumetría Operacional & 340 tractocamiones; $\approx 88.000$ viajes/año & 310 unidades de frío; $\approx 82.000$ viajes/año & 390 tractocamiones; $\approx 98.000$ viajes/año & Comparable con 374 camiones y 96.000 viajes/año \\
Arquitectura Tecnológica & Híbrida (Borde en cabina + Azure Cloud) & Híbrida (Edge on-premise + Azure Cloud) & Híbrida (Nube Azure + Servidores locales) & Al menos 1 proyecto con arquitectura híbrida \\
Nivel de Servicio (SLA) & 99,5\% mensual de disponibilidad & 99,2\% mensual; 72 h tolerancia offline & 99,6\% mensual de disponibilidad & Curimón: disponibilidad crítica $\ge 99{,}9\%$ \\
Estado y Ejecución & Finalizado; 100\% audIT como Principal & Finalizado; 100\% audIT como Principal & Finalizado; 100\% audIT como Principal & Concluidos y en operación últimos 5 años \\
\end{tabla}

El análisis de la \tabref{1-proyectos} evidencia que audIT ha resuelto con anterioridad y éxito comprobado los desafíos medulares que enfrenta Transportes Curimón S.A.:

\begin{enumerate}
  \item El Proyecto 1 acredita la capacidad de capturar, normalizar y transmitir eventos telemáticos masivos sobre flotas interurbanas de gran envergadura (340 tractos) operando en el mismo corredor geográfico de la Ruta 5, integrando lectura no intrusiva de bus CAN J1939 y geocercas operacionales.
  \item El Proyecto 2 describe una operación con conectividad intermitente, almacenamiento local y sincronización posterior, con una autonomía declarada de 72 horas y monitoreo de 310 unidades refrigeradas. Estos antecedentes requieren respaldo del Formulario T-6.
  \item El Proyecto 3 ratifica la capacidad de orquestar arquitecturas cloud híbridas de alto desempeño en Microsoft Azure capaces de ingerir más de 12 millones de transacciones diarias con una disponibilidad auditada del 99,6\% mensual.
\end{enumerate}

El detalle exhaustivo de los proyectos, el desglose de los 11 campos reglamentarios y los datos de contacto de las contrapartes técnicas que certifican la veracidad de estos antecedentes se presentan de forma pormenorizada en el anexo independiente AUDIT-Formulario-T-6.pdf, conforme a lo ordenado por las Bases Administrativas y el Comunicado 10.

\subsection{Certificaciones institucionales y alineamiento normativo}\label{sec:1-certificaciones}

audIT sustenta su práctica de ingeniería en un sólido marco de certificaciones corporativas y cumplimiento de estándares internacionales y normativas chilenas:

\begin{itemize}
  \item \textbf{ISO 9001:2015 (Sistema de Gestión de la Calidad):} Certificación corporativa emitida y vigente para el diseño, desarrollo, pruebas, implantación, integración de hardware y soporte continuo de plataformas de software y sistemas IoT telemáticos.
  \item \textbf{ISO/IEC 27001:2022 (Sistema de Gestión de Seguridad de la Información):} El proceso de certificación se encuentra en auditoría Fase 2 con Bureau Veritas. Conforme al mecanismo del \art{34.1}{24}, se compromete la entrega del certificado formal durante el Mes 1 del contrato, según el plan institucional incorporado en el Anexo 1.A.
  \item \textbf{Estándares Técnicos Complementarios de Ingeniería (\art{4.3}{5}):}
  \begin{itemize}
    \item \textit{NIST SP 800-207:} Arquitectura de Confianza Cero (\textit{Zero Trust}) como referencia para diseñar controles de acceso a recursos; la aplicación concreta y su evidencia se validarán durante la arquitectura de solución (\parencite{nist800207}).
    \item \textit{NIST Cybersecurity Framework 2.0 (CSF 2.0):} Marco de referencia para organizar resultados de gestión de ciberseguridad en torno a gobernanza, identificación, protección, detección, respuesta y recuperación (\parencite{nistcsf20}).
    \item \textit{OWASP ASVS 4.0 Nivel 2 y OWASP Top 10:} Guías mandatorias de desarrollo seguro aplicadas a la totalidad de microservicios y portales web (\parencite{owasp_asvs4}).
    \item \textit{SLSA Nivel 3 (Supply-chain Levels for Software Artifacts):} Trazabilidad criptográfica de compilaciones y pipelines de CI/CD para blindar la cadena de suministro de software (\caso{RT-11.28}{33}; \parencite{slsav1}).
    \item \textit{ISO/IEC 20000-1 e ISO 22301:} Referencias para los procesos de gestión de servicios y continuidad operacional; la conformidad se limita a ISO 9001:2015 vigente y no se declara certificación ni conformidad acreditada con estas normas adicionales (\parencite{iso20000,iso22301}).
  \end{itemize}
  \item \textbf{Competencias Regulatorias y Sectoriales Acreditadas:}
  \begin{itemize}
    \item \textit{Artículo 25 bis del Código del Trabajo:} La revisión del régimen aplicable considera, entre otros aspectos, el límite de cinco horas de conducción continua, los descansos correspondientes y la jornada mensual vigente a la fecha de la oferta. Los criterios de control se deben validar con asesoría competente antes de configurarse en el sistema (\parencite{codigo_trabajo_25bis,res_ex_1213_2009}).
    \item \textit{Decreto Supremo N.° 298/1994 (MTT):} Reglamento de transporte terrestre de cargas peligrosas en ruta (\parencite{ds298}).
    \item \textit{Decreto Supremo N.° 43/2015 (MINSAL):} Reglamento de almacenamiento de sustancias peligrosas en instalaciones industriales y patios químicos (\parencite{ds43minsal}).
    \item \textit{Protección de datos personales:} La fecha de oferta consignada en S2 es 5 de octubre de 2026, anterior a la entrada en vigor de la Ley N.° 21.719 prevista para el 1 de diciembre de 2026. Por ello, debe confirmarse si la entrega de la oferta se realizará después de esa fecha y actualizar el marco aplicable; rige la Ley N.° 19.628 hasta la entrada en vigor de la reforma (\parencite{ley19628,ley21719}).
    \item \textit{Ley N.° 21.663 (Ley Marco de Ciberseguridad):} Adecuación a las obligaciones de notificación y resguardo de infraestructura crítica de la información (\parencite{ley21663}).
    \item \textit{Ley N.° 21.377 (Ley No Chat):} Enclavamiento cinético vehicular por software y hardware para inhibir pantallas táctiles en cabina durante la marcha (\parencite{ley21377}).
  \end{itemize}
\end{itemize}

Las referencias oficiales a la legislación, los marcos NIST y las normas ISO citadas en este subdocumento se incorporan a la bibliografía compartida del paquete.

\subsection{Antecedentes financieros}\label{sec:1-financiero}

Los ratios financieros, estados de respaldo y antecedentes de solvencia corporativa se presentan en el Sobre N.º 1. Esta sección técnica no reproduce esos indicadores y remite al expediente financiero correspondiente.

Las tarifas, honorarios y valores monetarios de la oferta económica se presentan exclusivamente en el Sobre Económico N.º 3, conforme al \art{50.2}{32}.

\section{Estructura para Proyecto}\label{sec:1-estructura-proyecto}

Para abordar el proyecto de Transporte de Carga para Transportes Curimón S.A., audIT adopta una estructura de proyecto dedicada y orientada a objetivos, concebida para maximizar la sincronización entre el desarrollo de software y el despliegue físico de campo.

La organización interna para el proyecto se estructura en cuatro frentes de trabajo funcionales bajo la conducción unificada de la Dirección de Proyecto:

\begin{enumerate}
  \item \textbf{Frente de Gestión Contractual y PMO Dedicada:} Liderado por el Jefe de Proyecto e integrado por analistas de la Oficina PMO. Administra la relación contractual directa con la contraparte técnica de Transportes Curimón S.A., gestiona la matriz de riesgos del proyecto, supervisa el cumplimiento de la Carta Gantt y coordina la entrega formal de los hitos de pago asociados a entregables.
  \item \textbf{Célula Alfa -- Plataforma Cloud, Backend y Capa Anticorrupción:} Conducida por el Arquitecto de Solución y el Líder de Datos. Asume la responsabilidad sobre el desarrollo del núcleo transaccional en Azure, la ingestión de eventos telemáticos, la base de datos distribuida, los motores de despacho algorítmico y la implementación de las interfaces API seguras para absorber la telemetría de los 192 camiones subcontratados con GPS preexistente.
  \item \textbf{Célula Beta -- Firmware Embarcado, Telemetría y Despliegue de Terreno:} Conducida por el Líder de Operación (Jefe de IoT y Terreno). Responsable de la adaptación del producto \textit{audIT EdgeHub v2.4}, el suministro e instalación física de pasarelas y pinzas inductivas \textit{CANclick} sobre los 87 camiones propios sin telemetría previa y los 34 camiones subcontratados sin GPS previo (totalizando 121 nuevas pasarelas embarcadas), la homologación y acople no intrusivo \textit{CANclick} sobre los 61 camiones propios con J1939 de fábrica (completando la intervención de hardware sobre 182 tractocamiones que, sumados a los 192 integrados por API en Célula Alfa, cubren el 100\% de los 374 tractocamiones de la flota), la sensorización térmica de las 44 ramplas de frío y la articulación técnica con la red de soporte en la Ruta 5.
  \item \textbf{Célula Gamma -- Calidad, Ciberseguridad y Verificación:} Liderada por los responsables propuestos de Calidad y Seguridad de la Información. Su alcance preliminar considera compuertas de calidad, métricas de pruebas, controles de integridad de registros de jornada y protección de datos. Los umbrales y mecanismos deben definirse en el plan de calidad; las obligaciones de privacidad deben revisarse conforme a la normativa vigente y a la entrada en vigor de la Ley N.° 21.719.
\end{enumerate}

La asignación de dedicación de cada frente de trabajo, la articulación de las células técnicas y los protocolos de coordinación operacional garantizan la cobertura integral de los requerimientos de implantación y operación 24/7 de Transportes Curimón S.A., operando con plena autosuficiencia técnica desde el hito de inicio de servicios.

\section{Alianzas}\label{sec:1-alianzas}

audIT complementa sus capacidades internas de ingeniería mediante alianzas tecnológicas institucionales consolidadas y convenios de soporte vigentes, las cuales transfieren valor técnico y respaldo directo a la operación de Transportes Curimón S.A.

A continuación se exponen las alianzas estratégicas vigentes en materia de infraestructura en la nube, hardware industrial y soporte operativo territorial:

\subsection{Alianza tecnológica de nube: Microsoft Azure}\label{sec:1-alianza-azure}

audIT ostenta la categoría de \textbf{Microsoft Cloud Solution Provider (CSP) Tier 1 (Direct Partner)}, con acreditación avanzada en infraestructura híbrida, tecnologías de Internet de las Cosas (\textit{Azure IoT Hub}, \textit{Azure Arc}) y bases de datos gestionadas de alta escala. Esta alianza institucional proporciona:

\begin{itemize}
  \item Acceso preferente a soporte técnico de ingeniería de Microsoft de nivel Premier (\textit{24/7 Enterprise Support}) con tiempos de escalamiento a especialistas de producto inferiores a treinta (30) minutos ante incidencias de plataforma en la nube.
  \item Créditos de arquitectura y revisión técnica de diseño (\textit{Well-Architected Framework Review}) con arquitectos de soluciones de Microsoft para validar la resiliencia multirregión y la optimización de consumo de recursos (\textit{FinOps}).
  \item Acceso anticipado a programas de seguridad avanzada, módulos HSM en la nube y herramientas de cumplimiento normativo global.
\end{itemize}

\subsection{Alianzas de hardware automotriz grado industrial e IoT}\label{sec:1-alianza-hardware}

La compañía mantiene acuerdos corporativos de aprovisionamiento directo y soporte de ingeniería con fabricantes y distribuidores autorizados de componentes electrónicos de grado automotriz:

\begin{itemize}
  \item \textbf{Proveedores de Semiconductores y Microelectrónica Industrial:} Convenios de soporte técnico directo con distribuidores oficiales de procesadores de grado automotriz (NXP Semiconductors / STMicroelectronics) y módulos de comunicación celular eIoT con certificación de antenas y compatibilidad de bandas chilenas B2/B4/B7/B28 (Quectel Wireless Solutions).
  \item \textbf{Fabricantes de Acopladores Inductivos FMS/CANbus:} Alianza con distribuidores de tecnología \textit{CANclick / Car-Clearance}, garantizando la disponibilidad de acopladores electromagnéticos certificados bajo directivas europeas de seguridad automotriz (UNECE R10), los cuales leen inductivamente el campo electromagnético del par trenzado CAN bus sin cortar ni perforar el aislante de cobre del cableado original del vehículo.
  \item \textbf{Instrumentación y Sensores Térmicos:} Convenios de provisión con fabricantes de sondas de temperatura PT100 encapsuladas en acero inoxidable grado alimentario AISI 316 con certificados individuales de calibración trazables al Instituto Nacional de Normalización (INN).
\end{itemize}

\subsection{Red estratégica de centros de soporte en terreno en Ruta 5}\label{sec:1-alianza-red}

En articulación directa con los cuatro nodos operacionales descritos en la sección~\ref{sec:1-red-regional}, audIT mantiene convenios marco y Acuerdos de Nivel Operacional (OLA) que norman la gobernanza técnica de la asistencia en terreno:

\begin{itemize}
  \item \textbf{Gobernanza y Acuerdos de Nivel Operacional (OLA):} Estandarización de protocolos de intervención mecánica y eléctrica automotriz, auditados semestralmente conforme a ISO 9001 e ISO/IEC 20000-1 para certificar la calibración del instrumental técnico.
  \item \textbf{Integración al Sistema de Mantenimiento:} Los talleres aliados operan integrados a la plataforma audIT mediante el módulo web PWA de asistencia técnica, registrando en tiempo real cada reemplazo de componente, diagnóstico de falla y número de serie con firma digital.
  \item \textbf{Custodia y Trazabilidad de Stock Estratégico:} Mantenimiento de un inventario regulatorio de seguridad en cada nodo (pasarelas \textit{EdgeHub}, sondas PT100 y acopladores \textit{CANclick}), auditado con trazabilidad serializada para asegurar reposición inmediata en ruta.
\end{itemize}

La gobernanza contractual, los acuerdos OLA y los protocolos de auditoría de servicio de la red de asistencia en ruta se integran formalmente a la operación desde el Mes 1 de servicios.

\nocite{cisbenchmarks2024, ley19799, ley21377, ley21663, ley21719, iso9001, iso12207, iso20000_1, iso22301, iso27001, ds43minsal, ds298mtt, dfl1codigo_trabajo, nist_sp800_207, nist_csf2, owasp_asvs4}

